\clearpage
\VerdierPhysicalPageAnchor{0176}%
\begingroup\small\setlength{\parskip}{0pt}\setlength{\emergencystretch}{1.5em}
\setlength{\abovedisplayskip}{4pt plus 2pt minus 2pt}%
\setlength{\belowdisplayskip}{4pt plus 2pt minus 2pt}%

\VerdierSemanticAnchor{ch3d-1-2}
\subsection{The Notable Subcategories of Derived Categories.}

\VerdierSemanticAnchor{ch3d-1-2-1}
\noindent\textbf{Proposition 1.2.1.}\quad
\textit{a) The subcategories
$K^{,*}(\mathcal A)\subset K(\mathcal A)$ are
$\operatorname{Ac}(\mathcal A)$-localizing on both the right and the left
(Chap. II, 2.3.5).}

\noindent\textit{b) The subcategories
$K^{+,*}(\mathcal A)\subset K(\mathcal A)$ are
$\operatorname{Ac}(\mathcal A)$-localizing on the right.}

\noindent\textit{b)\textsuperscript{bis} The subcategories
$K^{-,*}(\mathcal A)\subset K(\mathcal A)$ are
$\operatorname{Ac}(\mathcal A)$-localizing on the left.}

\noindent\textit{c) The subcategory
$K^{b}(\mathcal A)\subset K^{+}(\mathcal A)$ is
$\operatorname{Ac}^{+}(\mathcal A)$-localizing on the left.}

\noindent\textit{c)\textsuperscript{bis} The subcategory
$K^{b}(\mathcal A)\subset K^{-}(\mathcal A)$ is
$\operatorname{Ac}^{-}(\mathcal A)$-localizing on the right.}

\noindent\textit{d) Every object of $K^{,+}(\mathcal A)$ admits a
quasi-isomorphism to an object of $K^{+}(\mathcal A)$.}

\noindent\textit{d)\textsuperscript{bis} Every object of
$K^{,-}(\mathcal A)$ admits a quasi-isomorphism from an object of
$K^{-}(\mathcal A)$.}

\noindent\textit{e) Every object of $K^{,b}(\mathcal A)$ admits a
quasi-isomorphism to an object of $K^{+,b}(\mathcal A)$.}

\noindent\textit{e)\textsuperscript{bis} Every object of
$K^{,b}(\mathcal A)$ admits a quasi-isomorphism from an object of
$K^{-,b}(\mathcal A)$.}

\noindent\textit{f) Every object of $K^{+,b}(\mathcal A)$ admits a
quasi-isomorphism from an object of $K^{b}(\mathcal A)$.}

\noindent\textit{f)\textsuperscript{bis} Every object of
$K^{-,b}(\mathcal A)$ admits a quasi-isomorphism to an object of
$K^{b}(\mathcal A)$.}

\VerdierSemanticAnchor{ch3d-1-2-1-truncations}
For every complex $X$, let $X(n,+\infty)$ (respectively, $X(-\infty,n)$)
denote the complex
\[
X(n,+\infty)=\cdots\longrightarrow0\longrightarrow0
\longrightarrow X^n/\operatorname{Im}(d_X^{n-1})
\longrightarrow X^{n+1}\longrightarrow\cdots
\]
(respectively,
\[
X(-\infty,n)=\cdots\longrightarrow X^{n-3}\longrightarrow X^{n-2}
\longrightarrow\operatorname{Ker}(d_X^{n-1})
\longrightarrow0\longrightarrow0\longrightarrow\cdots
\])
and let $p\colon X\longrightarrow X(n,+\infty)$ (respectively,
$i\colon X(-\infty,n)\longrightarrow X$) be the morphism of complexes
given by
\[
\begin{aligned}
p^q&=\operatorname{id},&&q>n,\\
p^n&\colon X^n\longrightarrow X^n/\operatorname{Im}(d_X^{n-1})
  &&\text{(canonical morphism)},\\
p^q&=0,&&q\leq n-1.
\end{aligned}
\]
\endgroup
